\chapter{Introduzione}

\section{Acustica}

In acustica ci interessa sia studiare cosa succede che effettivamente produrre
suoni e vedere che succede\footnote{Questa frase a lezione era un dissing ad
astronomia}. In generale l'acustica si occupa di generazione, propagazione e
ricezione di onde sonore o acustiche. È una \textbf{perturbazione}
\textbf{meccanica} che si \textbf{propaga} in un \textbf{fluido}. Perturbazione
perché riguarda piccoli spostamenti, non è uno shock. Meccanica perché riguarda
i movimenti delle particelle. Si propaga significa che è coinvolto qualcosa che
cambia a seconda dell'istante di tempo in punti spaziali diversi. Infine in un
fluido perché guardiamo onde in un gas o liquido.

Per derivare la funzione d'onda, immaginiamo di avere una funzione \(\psi
{(\mathbf{x} , t)}\) che è una densità che dipende sia dal tempo che dallo
spazio. Ci interessa come viene modificata questa densità nel tempo, quindi
guardiamo l'equazione
\begin{equation}\label{eq:acu_onda_integrale}
  \frac{\partial}{\partial t}\int _D \psi {(\mathbf{x} , t)} \,d\mathbf{x} =
  \int _D q{(\mathbf{x} , t)} \,d\mathbf{x} - \oint _{\partial D} \mathbf{J}
  \cdot \mathbf{n} \,dS
\end{equation}
dove \(q\) è il termine di \emph{sorgente}, che dice quanta di questa densità
viene generata nell'istante \(t\) al punto \(\mathbf{x} \), \(\mathbf{n} \) è il
vettore normale alla superficie \(S\) e \(\mathbf{J} \) è il flusso. Vogliamo
inoltre che il flusso \(\mathbf{J} = \psi \mathbf{v} \). Usando dunque ora
il teorema della divergenza otteniamo che
\begin{equation}\label{eq:acu_div}
  \int _D {\left( \frac{\partial }{\partial t}\psi - q + \mathrm{div} {(\psi \mathbf{v}
  )}\right)}  \,d\mathbf{x} = 0 \implies \frac{\partial }{\partial t}\psi +
  \mathrm{div} {(\psi \mathbf{v})} = q
\end{equation}

In questo caso sappiamo che \(\psi = \rho\) la densità di massa, e poiché
sappiamo che non è possibile creare massa, allora \(q = 0\). otteniamo dunque
l'equazione
\begin{equation}\label{eq:acu_cont_massa}
  \frac{\partial \rho}{\partial t} + \mathrm{div}{(\rho \mathbf{v} )} = 0
\end{equation}

Un'altra ``densità'' che ci interessa è la \emph{quantità di moto}. Possiamo
infatti prendere dunque \(\psi = {(\rho \mathbf{v} )}_j\), per \(j=1,2,3\). Se fosse infatti che
\(\mathbf{v} \) è costante, avremmo che \(\int \rho \mathbf{v}  = \mathbf{v}
\int _D \rho = \mathbf{v} m\). Per considerare l'inerzia ora includiamo la legge
di Newton \(\mathbf{F} = {(m\mathbf{v} )}'\), allora abbiamo che \(q = F_j\).
Riscrivendo dunque l'equazione~\eqref{eq:acu_div}, abbiamo
\[
  \frac{\partial }{\partial t}{(\rho v_j)} + \mathrm{div}{(\rho v_s \mathbf{v}
  )} = F_j = {(\nabla p)}_j
\]
con \(p\) la pressione. Otteniamo ora, con la derivata del prodotto
\[
  v_j {\left( \frac{\partial \rho}{\partial t} +  \mathrm{div}{(\rho \mathbf{v}
  )}\right)} + \rho \frac{\partial v_j}{\partial t} + \rho \mathbf{v} \cdot
  \nabla v_j = {(\nabla p)}_j
\]
ma da~\eqref{eq:acu_cont_massa} abbiamo che il termine tra parentesi è nullo, da
cui (riscrivendo il tutto anche in forma vettoriale e dividendo per \(\rho\) )
\begin{equation}\label{eq:eulero}
  \frac{\partial \mathbf{v} }{\partial t} + {(\mathbf{v} \cdot \nabla )}
  \mathbf{v} = \frac{1}{\rho} \nabla p
\end{equation}
dove \({\left( {(\mathbf{v} \cdot \nabla )}\mathbf{v}  \right)}_j =
\sum_{k=1}^{3} v_{n} \frac{\partial v_j}{\partial x_k}  \) 

